% Môn: Vật lý
% Lớp: 10
% Chương: Chương I. Mở đầu
% Bài: L10C1 Bài 2. Các quy tắc an toàn trong phòng thực hành Vật lí · Giải nghĩa các kí hiệu kĩ thuật và biển cảnh báo an toàn trong phòng thực hành.
% Số câu: 185
% App đọc file này trực tiếp — sửa rồi Commit trên GitHub, bấm Đồng bộ GitHub .tex

% L10C1 Bài 2. Các quy tắc an toàn trong phòng thực hành Vật lí



% ===== Câu 75 =====
% ID: L10C1B2-16-DS
% Mức: VD
\dangbt{Giải nghĩa các kí hiệu kĩ thuật và biển cảnh báo an toàn trong phòng thực hành.}
\begin{ex}
% ID: L10C1B2-16-DS
% Mức: VD
Một học sinh chuẩn bị thực hành thí nghiệm điện trong phòng thực hành Vật lí. Em quan sát thấy bốn kí hiệu/biển cảnh báo được gắn tại các vị trí khác nhau trong phòng và ghi lại nhận xét như sau. Hãy xét tính đúng/sai của từng nhận định.
\choiceTF
{\True Biển cảnh báo hình tam giác viền vàng, bên trong có hình tia sét màu đen cảnh báo nguy hiểm về điện — khu vực hoặc thiết bị có điện áp cao, không được chạm tay trực tiếp vào.}
{Biển cấm hình tròn viền đỏ, bên trong có hình ngọn lửa bị gạch chéo có nghĩa là cấm sử dụng nước trong khu vực đó để tránh trơn trượt.}
{\True Kí hiệu hình đầu lâu xương chéo màu đen trên nền trắng (hoặc vàng) cảnh báo chất độc hại — học sinh không được tự ý mở, ngửi hoặc tiếp xúc trực tiếp với hóa chất được dán nhãn này.}
{Biển chỉ dẫn hình vuông nền xanh lá, bên trong có hình người chạy về phía cửa thoát hiểm là biển cảnh báo khu vực nguy hiểm, học sinh phải tránh xa khu vực đó trong suốt buổi thực hành.}
\loigiai{
\begin{itemchoice}
\item (Đúng) Biển cảnh báo hình tam giác viền vàng, bên trong có hình tia sét màu đen cảnh báo nguy hiểm về điện — khu vực hoặc thiết bị có điện áp cao, không được chạm tay trực tiếp vào. (Vì): Biển tam giác viền vàng có hình tia sét là kí hiệu cảnh báo nguy hiểm điện áp cao, yêu cầu không chạm tay trực tiếp vào thiết bị hoặc khu vực đó
\item (Sai) Biển cấm hình tròn viền đỏ, bên trong có hình ngọn lửa bị gạch chéo có nghĩa là cấm sử dụng nước trong khu vực đó để tránh trơn trượt. (Vì): Biển cấm hình tròn viền đỏ có hình ngọn lửa bị gạch chéo có nghĩa là cấm lửa/cấm chất dễ cháy (no open flame), không phải cấm sử dụng nước
\item (Đúng) Kí hiệu hình đầu lâu xương chéo màu đen trên nền trắng (hoặc vàng) cảnh báo chất độc hại — học sinh không được tự ý mở, ngửi hoặc tiếp xúc trực tiếp với hóa chất được dán nhãn này. (Vì): Kí hiệu đầu lâu xương chéo là biểu tượng quốc tế cảnh báo chất độc hại; học sinh tuyệt đối không tự ý tiếp xúc, ngửi hoặc mở các hóa chất mang nhãn này
\item (Sai) Biển chỉ dẫn hình vuông nền xanh lá, bên trong có hình người chạy về phía cửa thoát hiểm là biển cảnh báo khu vực nguy hiểm, học sinh phải tránh xa khu vực đó trong suốt buổi thực hành. (Vì): Biển hình vuông nền xanh lá có hình người chạy về phía cửa là biển chỉ dẫn lối thoát hiểm khẩn cấp, không phải biển cảnh báo khu vực nguy hiểm cần tránh xa
\end{itemchoice}
}
\end{ex}

% ===== Câu 78 =====
% ID: L10C1B2-17-DS
% Mức: VD
\begin{ex}
% ID: L10C1B2-17-DS
% Mức: VD
Trong phòng thực hành Vật lí, học sinh quan sát thấy bốn kí hiệu/biển cảnh báo sau đây được dán trên các thiết bị và khu vực khác nhau. Hãy xét tính đúng/sai của các nhận định về ý nghĩa của từng kí hiệu đó.
\choiceTF
{\True Biển hình tam giác viền vàng, bên trong có hình tia sét màu đen cảnh báo nguy hiểm về điện – khu vực hoặc thiết bị có điện áp cao, có thể gây điện giật chết người nếu chạm vào.}
{Biển hình tròn có nền xanh lam, bên trong có hình người đang đeo kính bảo hộ có nghĩa là 'Cấm đeo kính bảo hộ trong khu vực này' nhằm tránh làm hỏng thiết bị quang học.}
{\True Biển hình tam giác viền đỏ, bên trong có hình ngọn lửa cảnh báo vật liệu hoặc chất dễ cháy – cần tránh xa nguồn nhiệt, tia lửa và không hút thuốc trong khu vực đó.}
{\True Biển hình tròn có nền đỏ, bên trong có hình bàn tay bị gạch chéo có nghĩa là 'Không được chạm tay vào' – thường dán trên các thiết bị đang hoạt động hoặc bề mặt có nhiệt độ cao.}
\loigiai{
\begin{itemchoice}
\item (Đúng) Biển hình tam giác viền vàng, bên trong có hình tia sét màu đen cảnh báo nguy hiểm về điện – khu vực hoặc thiết bị có điện áp cao, có thể gây điện giật chết người nếu chạm vào. (Vì): Biển tam giác viền vàng với hình tia sét là kí hiệu cảnh báo nguy hiểm điện (điện áp cao), cảnh báo nguy cơ điện giật nghiêm trọng hoặc tử vong nếu tiếp xúc trực tiếp
\item (Sai) Biển hình tròn có nền xanh lam, bên trong có hình người đang đeo kính bảo hộ có nghĩa là 'Cấm đeo kính bảo hộ trong khu vực này' nhằm tránh làm hỏng thiết bị quang học. (Vì): Biển hình tròn nền xanh lam với hình người đeo kính bảo hộ là biển BẮ $T$ BUỘ $C$ đeo kính bảo hộ (lệnh bắt buộc), không phải biển cấm; mục đích là bảo vệ mắt người dùng khỏi hóa chất bắn hoặc mảnh vỡ
\item (Đúng) Biển hình tam giác viền đỏ, bên trong có hình ngọn lửa cảnh báo vật liệu hoặc chất dễ cháy – cần tránh xa nguồn nhiệt, tia lửa và không hút thuốc trong khu vực đó. (Vì): Biển tam giác viền đỏ (hoặc vàng) với hình ngọn lửa là kí hiệu cảnh báo chất/vật liệu dễ cháy, yêu cầu tránh xa nguồn nhiệt, tia lửa và nghiêm cấm hút thuốc trong khu vực đó
\item (Đúng) Biển hình tròn có nền đỏ, bên trong có hình bàn tay bị gạch chéo có nghĩa là 'Không được chạm tay vào' – thường dán trên các thiết bị đang hoạt động hoặc bề mặt có nhiệt độ cao. (Vì): Biển hình tròn nền đỏ với hình bàn tay bị gạch chéo là biển cấm chạm tay vào, thường được dán trên thiết bị đang vận hành, bề mặt nóng hoặc bộ phận chuyển động nguy hiểm
\end{itemchoice}
}
\end{ex}

% ===== Câu 81 =====
% ID: L10C1B2-18-DS
% Mức: VD
\begin{ex}
% ID: L10C1B2-18-DS
% Mức: VD
Trong buổi thực hành Vật lí, giáo viên yêu cầu học sinh đọc và hiểu ý nghĩa của bốn kí hiệu/biển cảnh báo được dán tại các vị trí khác nhau trong phòng thí nghiệm trước khi tiến hành thí nghiệm. Dưới đây là bốn nhận định về các kí hiệu đó. Hãy xác định mỗi nhận định là Đúng hay Sai.
\choiceTF
{Biển cảnh báo hình tam giác viền vàng, bên trong có hình tia sét màu đen, dán trên tủ điện hoặc ổ cắm điện, có nghĩa là khu vực đó an toàn, học sinh có thể tự ý cắm thiết bị vào nguồn điện mà không cần hỏi giáo viên.}
{\True Kí hiệu hình ngọn lửa màu đỏ/cam được dán trên chai hóa chất hoặc bình gas trong phòng thực hành cảnh báo chất dễ cháy, học sinh cần tránh xa nguồn nhiệt và không để gần lửa khi sử dụng.}
{Biển cảnh báo hình đầu lâu xương chéo (skull and crossbones) dán trên bình chứa hóa chất có nghĩa là hóa chất bên trong vô hại, học sinh có thể tự do lấy và sử dụng mà không cần trang bị bảo hộ.}
{\True Kí hiệu hình mắt kính bảo hộ (goggles) được dán ở lối vào khu vực thí nghiệm có nghĩa là bắt buộc đeo kính bảo hộ khi vào khu vực đó để bảo vệ mắt khỏi các mảnh vỡ, hóa chất bắn hoặc tia sáng nguy hiểm.}
\loigiai{
\begin{itemchoice}
\item (Sai) Biển cảnh báo hình tam giác viền vàng, bên trong có hình tia sét màu đen, dán trên tủ điện hoặc ổ cắm điện, có nghĩa là khu vực đó an toàn, học sinh có thể tự ý cắm thiết bị vào nguồn điện mà không cần hỏi giáo viên. (Vì): Biển cảnh báo hình tam giác viền vàng có hình tia sét màu đen là kí hiệu nguy hiểm điện áp cao, cảnh báo nguy cơ giật điện; học sinh tuyệt đối không được tự ý thao tác với nguồn điện mà không có sự cho phép và hướng dẫn của giáo viên
\item (Đúng) Kí hiệu hình ngọn lửa màu đỏ/cam được dán trên chai hóa chất hoặc bình gas trong phòng thực hành cảnh báo chất dễ cháy, học sinh cần tránh xa nguồn nhiệt và không để gần lửa khi sử dụng. (Vì): Kí hiệu hình ngọn lửa cảnh báo chất dễ cháy (flammable); khi làm việc với các chất này, học sinh phải giữ xa nguồn nhiệt, tia lửa và ngọn lửa hở để tránh nguy cơ cháy nổ
\item (Sai) Biển cảnh báo hình đầu lâu xương chéo (skull and crossbones) dán trên bình chứa hóa chất có nghĩa là hóa chất bên trong vô hại, học sinh có thể tự do lấy và sử dụng mà không cần trang bị bảo hộ. (Vì): Biển hình đầu lâu xương chéo là kí hiệu quốc tế cảnh báo chất độc hại hoặc gây chết người; học sinh phải mặc đồ bảo hộ, đeo găng tay và khẩu trang, tuyệt đối không tự ý tiếp xúc hay sử dụng hóa chất đó
\item (Đúng) Kí hiệu hình mắt kính bảo hộ (goggles) được dán ở lối vào khu vực thí nghiệm có nghĩa là bắt buộc đeo kính bảo hộ khi vào khu vực đó để bảo vệ mắt khỏi các mảnh vỡ, hóa chất bắn hoặc tia sáng nguy hiểm. (Vì): Kí hiệu hình kính bảo hộ là biển bắt buộc (mandatory sign) yêu cầu đeo kính bảo hộ trong khu vực đó nhằm bảo vệ mắt khỏi các tác nhân nguy hiểm như mảnh vỡ thủy tinh, hóa chất bắn hoặc bức xạ
\end{itemchoice}
}
\end{ex}

% ===== Câu 122 =====
% ID: L10C1B2-37-TN
% Mức: NB
\begin{ex}
% ID: L10C1B2-37-TN
% Mức: NB
Biển cảnh báo có hình tam giác đều viền đen, nền vàng, bên trong vẽ biểu tượng một tia sét màu đen chỉ xuống dưới dùng để cảnh báo mối nguy hiểm nào sau đây?
\choice
{Nguy hiểm về nhiệt độ cao}
{Nguy hiểm về từ trường mạnh}
{\True Nguy hiểm về điện}
{Nguy hiểm về chất phóng xạ}
\loigiai{
Biểu tượng tia sét màu đen đặt trong hình tam giác viền đen nền vàng là kí hiệu thuộc nhóm cảnh báo nguy hiểm, cụ thể là cảnh báo nguy hiểm về điện.
}
\end{ex}

% ===== Câu 123 =====
% ID: L10C1B2-38-TN
% Mức: NB
\begin{ex}
% ID: L10C1B2-38-TN
% Mức: NB
Kí hiệu "+" hoặc màu đỏ mang ý nghĩa là
\choice
{đầu vào}
{đầu ra}
{\True cực dương}
{cực âm}
\loigiai{
Kí hiệu "+" hoặc màu đỏ thường dùng để chỉ cực dương.
}
\end{ex}

% ===== Câu 126 =====
% ID: L10C1B2-39-TN
% Mức: NB
\begin{ex}
% ID: L10C1B2-39-TN
% Mức: NB
Biển báo có biểu tượng \textbf{hai ống nghiệm đang nhỏ chất lỏng xuống làm mòn mặt phẳng và bàn tay} (đặt trong hình tam giác vàng) cảnh báo điều gì?
\choice
{\True Cảnh báo hóa chất ăn mòn.}
{Nước dùng để rửa tay.}
{Chất độc môi trường.}
{Cảnh báo vật sắc nhọn.}
\loigiai{
Biểu tượng này cảnh báo khu vực hoặc bình chứa hóa chất có tính ăn mòn cao, cần cẩn thận khi tiếp xúc.
}
\end{ex}

% ===== Câu 130 =====
% ID: L10C1B2-41-TN
% Mức: NB
\begin{ex}
% ID: L10C1B2-41-TN
% Mức: NB
Biển báo có hình chữ nhật (hoặc hình vuông), nền đỏ, bên trong có biểu tượng \textbf{bình chữa cháy} màu trắng là loại biển báo gì?
\choice
{Biển cảnh báo cháy nổ.}
{Biển cấm lửa.}
{\True Biển thông báo vị trí phương tiện chữa cháy.}
{Biển báo cấm sử dụng bình chữa cháy.}
\loigiai{
Nhóm biển hình vuông/chữ nhật mang ý nghĩa thông báo, chỉ dẫn vị trí đặt các thiết bị an toàn (như bình chữa cháy) hoặc lối thoát hiểm.
}
\end{ex}

% ===== Câu 132 =====
% ID: L10C1B2-42-TN
% Mức: TH
\begin{ex}
% ID: L10C1B2-42-TN
% Mức: TH
Khi quan sát vỏ hộp chứa một thiết bị đo quang học nhạy cảm, học sinh thấy có vẽ biểu tượng một chiếc ô màu đen kèm các tia sáng thẳng đứng chiếu xuống phía trên chiếc ô. Kí hiệu này nhắc nhở người bảo quản cần lưu ý điều gì?
\choice
{Tránh để thiết bị tiếp xúc trực tiếp với nước mưa}
{\True Tránh ánh nắng mặt trời chiếu trực tiếp vào thiết bị}
{Phải đặt đứng thiết bị theo chiều mũi tên hướng lên}
{Không được để thiết bị ở những nơi có độ ẩm cao}
\loigiai{
Kí hiệu chiếc ô che tia sáng thẳng đứng biểu thị yêu cầu bảo quản dụng cụ tránh ánh nắng chiếu trực tiếp từ mặt trời để bảo vệ thấu kính hoặc linh kiện bên trong.
}
\end{ex}

% ===== Câu 134 =====
% ID: L10C1B2-43-TN
% Mức: TH
\begin{ex}
% ID: L10C1B2-43-TN
% Mức: TH
Biển báo có hình tròn viền đỏ, nền trắng, có một vạch đỏ chéo đè lên hình vẽ ngọn lửa màu đen thuộc loại biển báo nào và có ý nghĩa gì?
\choice
{Biển cảnh báo nguy hiểm, có nghĩa là khu vực có chất dễ cháy}
{Biển chỉ dẫn, có nghĩa là vị trí đặt bình chữa cháy}
{\True Biển cấm, có nghĩa là cấm lửa}
{Biển báo bắt buộc thực hiện, có nghĩa là phải trang bị bảo hộ chống cháy}
\loigiai{
Biển có hình tròn viền đỏ có vạch chéo đỏ đè lên hình vẽ biểu tượng màu đen là biển báo cấm, hình ngọn lửa biểu thị cho lửa, do đó biển này có ý nghĩa cấm lửa.
}
\end{ex}

